%% Created by Maple 2022.1, Windows 10
%% Source Worksheet: lab_sols03
%% Generated: Wed May 08 10:31:51 BST 2024
\documentclass{article}
\usepackage{amssymb}
\usepackage{graphicx}
\usepackage{hyperref}
\usepackage{listings}
\usepackage{mathtools}
\usepackage{maple}
\usepackage[utf8]{inputenc}
\usepackage[svgnames]{xcolor}
\usepackage{amsmath}
\usepackage{breqn}
\usepackage{textcomp}
\begin{document}
\lstset{basicstyle=\ttfamily,breaklines=true,columns=flexible}
\pagestyle{empty}
\DefineParaStyle{Maple Bullet Item}
\DefineParaStyle{Maple Heading 1}
\DefineParaStyle{Maple Warning}
\DefineParaStyle{Maple Heading 4}
\DefineParaStyle{Maple Heading 2}
\DefineParaStyle{Maple Heading 3}
\DefineParaStyle{Maple Dash Item}
\DefineParaStyle{Maple Error}
\DefineParaStyle{Maple Title}
\DefineParaStyle{Maple Ordered List 1}
\DefineParaStyle{Maple Text Output}
\DefineParaStyle{Maple Ordered List 2}
\DefineParaStyle{Maple Ordered List 3}
\DefineParaStyle{Maple Normal}
\DefineParaStyle{Maple Ordered List 4}
\DefineParaStyle{Maple Ordered List 5}
\DefineCharStyle{Maple 2D Output}
\DefineCharStyle{Maple 2D Input}
\DefineCharStyle{Maple Maple Input}
\DefineCharStyle{Maple 2D Math}
\DefineCharStyle{Maple Hyperlink}
\begin{center}
\begin{Maple Title}
\underline{\textbf{Plotting}}
\end{Maple Title}
\end{center}
\begin{Maple Normal}
\_\_\_\_\_\_\_\_\_\_\_\_\_\_\_\_\_\_\_\_\_\_\_\_\_\_\_\_\_\_\_\_\_\_\_\_\_\_\_\_\_\_\_\_\_\_\_\_\_\_\_\_\_\_\_\_\_\_\_\_\_\_\_\_\_\_\_\_\_\_\_\_\_\_\_\_\_\_\_\_\_
\end{Maple Normal}
\begin{Maple Heading 2}
\textbf{Exercise 1.1}
\end{Maple Heading 2}
\begin{lstlisting}
> restart;
\end{lstlisting}
\begin{Maple Normal}
\textbf{(a)}
\end{Maple Normal}
\begin{lstlisting}
> plot(2*exp(-t)*sin(30*t),t=-1..5);
\end{lstlisting}
\begin{Maple Normal}
This oscillates and dies away, like the motion of a guitar string after it has been plucked.
\end{Maple Normal}
\begin{Maple Normal}
\textbf{(b)}
\end{Maple Normal}
\begin{lstlisting}
> plot(2*sin(20*t)+3*sin(21*t),t=0..30,numpoints=1000);
\end{lstlisting}
\begin{Maple Normal}
This oscillates rapidly (at the same frequency as 
$\sin ( 20.5 t)$), with the size of the oscillations varying more slowly  (like 
$\cos (t)$).
\end{Maple Normal}
\begin{Maple Normal}
\textbf{(c)}
\end{Maple Normal}
\begin{lstlisting}
> plot(sin(x)+sin(3*x)/3+sin(5*x)/5+sin(7*x)/7,x=0..20);
\end{lstlisting}
\begin{Maple Normal}
Half the time this wiggles around near 
$y = 0.8$, and half the time it wiggles around near 
$y =- 0.8$.  It jumps quite rapidly between the two levels, and the jumps occur at multiples of 
$\pi$ (
$\pi= 3.14,2 \pi= 6.28,3 \pi= 9.42$ and so on).
\end{Maple Normal}
\begin{Maple Normal}
\textbf{(d)}
\end{Maple Normal}
\begin{lstlisting}
> plot(cos(Pi*x^2),x=-5..5);
\end{lstlisting}
\begin{Maple Normal}
The function oscillates between 
$-1$ and 
$1$, with the oscillations becoming more and more rapid as the absolute value of 
$x$ increases.
\end{Maple Normal}
\begin{Maple Normal}
\_\_\_\_\_\_\_\_\_\_\_\_\_\_\_\_\_\_\_\_\_\_\_\_\_\_\_\_\_\_\_\_\_\_\_\_\_\_\_\_\_\_\_\_\_\_\_\_\_\_\_\_\_\_\_\_\_\_\_\_\_\_\_\_\_\_\_\_\_\_\_\_\_\_\_\_\_\_\_\_\_
\end{Maple Normal}
\begin{Maple Heading 2}
\textbf{Exercise 1.2}
\end{Maple Heading 2}
\begin{lstlisting}
> restart;
\end{lstlisting}
\begin{lstlisting}
> f := (x) -> (x^3-x)*(x^2-4/9)*(x^2-1/9);
\end{lstlisting}
\begin{Maple Normal}
The graph is quite flat and close to zero between 
$x =-1$ and 
$x =1$, crossing the 
$ x$-axis seven times, at 
$x =-1,-\frac{2}{3},-\frac{1}{3},0,\frac{1}{3},\frac{2}{3}$  and 
$1$.  The function increases rapidly towards 
$\infty$ for 
$1<x$, and drops rapidly towards 
$-\infty$ for 
$x <-1$.  
\end{Maple Normal}
\begin{lstlisting}
> plot(f(x),x=-2..2,-0.5..0.5);
\end{lstlisting}
\begin{Maple Normal}
The rate of growth is similar to 
$x^{7}$, and in fact 
$f (x)$ is indistinguishable from 
$x^{7}$ when 
$x$ is reasonably large.
\end{Maple Normal}
\begin{lstlisting}
> plot([f(x),x^7],x=-3..3);
\end{lstlisting}
\begin{lstlisting}
> plot([f(x),x^7],x=-10..10);
\end{lstlisting}
\begin{Maple Normal}
To see the reason for this, just expand out 
$f (x)$.  The leading term is 
$x^{7}$, and the remaining terms are multiples of 
$x^{5}$ or lower, so they are much smaller than 
$x^{7}$ when 
$x$ is large.
\end{Maple Normal}
\begin{lstlisting}
> expand(f(x));
\end{lstlisting}
\begin{Maple Normal}
\_\_\_\_\_\_\_\_\_\_\_\_\_\_\_\_\_\_\_\_\_\_\_\_\_\_\_\_\_\_\_\_\_\_\_\_\_\_\_\_\_\_\_\_\_\_\_\_\_\_\_\_\_\_\_\_\_\_\_\_\_\_\_\_\_\_\_\_\_\_\_\_\_\_\_\_\_\_\_\_\_
\end{Maple Normal}
\begin{Maple Heading 2}
\textbf{Exercise 1.3}
\end{Maple Heading 2}
\begin{lstlisting}
> restart;
\end{lstlisting}
\begin{lstlisting}
> g := (x) -> (2*x-3)/(3*x-4);
\end{lstlisting}
\begin{Maple Normal}
\textbf{(a)} Most of the time the graph is very flat, and close to 
$y = 0.7$.  Somewhere close to 
$x = 1.3$, the graph climbs steeply to 
$\infty$, jumps down discontinuously to 
$-\infty$, and then climbs back up to around 
$ 0.7$ again.
\end{Maple Normal}
\begin{lstlisting}
> plot(g(x),x=-10..10,-10..10,numpoints=1000,discont=true);
\end{lstlisting}
\begin{lstlisting}
> plot(g(x),x=-20..20,0..2,axes=boxed);
\end{lstlisting}
\begin{lstlisting}
> plot(g(x),x=1..2,-2..4);
\end{lstlisting}
\begin{lstlisting}
> 
\end{lstlisting}
\begin{Maple Normal}
\textbf{(b) }The function is discontinuous at the point where the formula 
$g (x)=\frac{2 x -3}{3 x -4}$ involves division by zero, which means 
$3 x -4=0$ or in other words 
$x =\frac{4}{3}$.   This can be found graphically as the point where the vertical line crosses the axis in the picture above.  Alternatively: 
\end{Maple Normal}
\begin{lstlisting}
> discont(g(x),x);
\end{lstlisting}
\begin{Maple Normal}
Here we replot the graph, skipping over the discontinuity:
\end{Maple Normal}
\begin{lstlisting}
>   plot(g(x),x=-3..5,-10..10,discont=true);
\end{lstlisting}
\begin{Maple Normal}
\textbf{(c)}
\end{Maple Normal}
\begin{lstlisting}
> plot([g(x),5],x=0..3,-10..10,discont=true);
\end{lstlisting}
\begin{Maple Normal}
This picture shows that the graph does not cross the line 
$y =\frac{2}{3}$
\end{Maple Normal}
\begin{lstlisting}
> plot([g(x),2/3],x=-20..20,0..1,discont=true);
\end{lstlisting}
\begin{Maple Normal}
To see this algebraically, we first find the general formula for the point where the curve crosses the line 
$y =a$:
\end{Maple Normal}
\begin{lstlisting}
> solve(g(x)=a,\{x\});
\end{lstlisting}
\begin{Maple Normal}
The formula involves division by zero when 
$a =\frac{2}{3}$, indicating that the curve does not in fact cross the line 
$y =\frac{2}{3}$.
\end{Maple Normal}
\begin{Maple Normal}
\textbf{(d)}
\end{Maple Normal}
\begin{lstlisting}
> limit(g(x),x=infinity);
\end{lstlisting}
\begin{lstlisting}
> limit(g(x),x=-infinity);
\end{lstlisting}
\begin{Maple Normal}
This is of course the same answer as in (c).  As 
$x$ approaches 
$\infty$, the function 
$g (x)$ approaches 
$\frac{2}{3}$ from below, but never reaches it.  As 
$x$ approaches 
$-\infty$, the function approaches 
$\frac{2}{3}$ from above, but again never reaches it.
\end{Maple Normal}
\begin{Maple Normal}
\_\_\_\_\_\_\_\_\_\_\_\_\_\_\_\_\_\_\_\_\_\_\_\_\_\_\_\_\_\_\_\_\_\_\_\_\_\_\_\_\_\_\_\_\_\_\_\_\_\_\_\_\_\_\_\_\_\_\_\_\_\_\_\_\_\_\_\_\_\_\_\_\_\_\_\_\_\_\_\_\_
\end{Maple Normal}
\begin{Maple Heading 2}
\textbf{Exercise 1.4}
\end{Maple Heading 2}
\begin{lstlisting}
> restart;
\end{lstlisting}
\begin{lstlisting}
> f := (x)->(1-exp(x)+exp(2*x)-exp(3*x))/sin(x);
\end{lstlisting}
\begin{lstlisting}
> f(0);
\end{lstlisting}
\begin{lstlisting}
> plot(f(x),x=-1..1);
\end{lstlisting}
\begin{Maple Normal}
Although the formula for 
$f (x)$ does not make sense when 
$x =0$, the only reasonable definition is 
$f (0)=-2$, as we see from the graph.  Here is a non-graphical way to get Maple to produce this answer:
\end{Maple Normal}
\begin{lstlisting}
> limit(f(x),x=0);
\end{lstlisting}
\begin{Maple Normal}
\_\_\_\_\_\_\_\_\_\_\_\_\_\_\_\_\_\_\_\_\_\_\_\_\_\_\_\_\_\_\_\_\_\_\_\_\_\_\_\_\_\_\_\_\_\_\_\_\_\_\_\_\_\_\_\_\_\_\_\_\_\_\_\_\_\_\_\_\_\_\_\_\_\_\_\_\_\_\_\_\_
\end{Maple Normal}
\begin{Maple Heading 2}
\textbf{Exercise 1.5}
\end{Maple Heading 2}
\begin{lstlisting}
> restart;
\end{lstlisting}
\begin{Maple Normal}
The Bessel function is even (ie 
$J_{2}(-x)=J_{2}(x)$).  It oscillates quite regularly, with amplitude decreasing slowly as 
$x$ approaches 
$\infty$ or 
$-\infty$.
\end{Maple Normal}
\begin{lstlisting}
> plot(BesselJ(2,x),x=-50..50);
\end{lstlisting}
\begin{Maple Normal}
For small 
$x$, the function 
$J_{2}(x)$ is very close to 
$\frac{x^{2}}{8}$.
\end{Maple Normal}
\begin{lstlisting}
> plot([BesselJ(2,x),x^2/5],x=-0.1..0.1);
\end{lstlisting}
\begin{lstlisting}
> plot([BesselJ(2,x),x^2/8],x=-0.1..0.1);
\end{lstlisting}
\begin{Maple Normal}
This plot has glitches because Maple has not calculated enough points.
\end{Maple Normal}
\begin{lstlisting}
> plot(BesselJ(2,x),x=10..300);
\end{lstlisting}
\begin{Maple Normal}
We can cure this with the \textcolor[RGB]{255,0,0}{\textbf{numpoints}} option.
\end{Maple Normal}
\begin{lstlisting}
> plot(BesselJ(2,x),x=10..300,numpoints=200);
\end{lstlisting}
\begin{lstlisting}
> pic := %:
\end{lstlisting}
\begin{Maple Normal}
Here are several attempts to find the 'envelope' of the graph.  The final attempt is 
$y = 0.8 x^{-\frac{1}{2}}$ , which fits very nicely with the peaks of the waves.
\end{Maple Normal}
\begin{lstlisting}
> with(plots):
\end{lstlisting}
\begin{lstlisting}
> display(pic,plot(6*x^(-1),x=10..300,colour=blue));
\end{lstlisting}
\begin{lstlisting}
> display(pic,plot(0.3*x^(-1/3),x=10..300,colour=blue));
\end{lstlisting}
\begin{lstlisting}
> display(pic,plot(0.5*x^(-1/2),x=10..300,colour=blue));
\end{lstlisting}
\begin{lstlisting}
> display(pic,plot(0.8*x^(-1/2),x=10..300,colour=blue));
\end{lstlisting}
\begin{Maple Normal}
This plot shows that 
$J_{2}(x)$ has (to a very good approximation) the same frequency as 
$\sin (x)$.
\end{Maple Normal}
\begin{lstlisting}
> plot([BesselJ(2,x),sin(x+Pi/4)],x=10..100);
\end{lstlisting}
\begin{Maple Normal}
By combining the above results, we see that the following function should be a good approximation to 
$J_{2}(x)$ for large 
$x$.  The plot below shows this graphically.
\end{Maple Normal}
\begin{lstlisting}
> f := (x) -> -4*sin(x+Pi/4)/(5*sqrt(x));
\end{lstlisting}
\begin{lstlisting}
> plot([BesselJ(2,x),f(x)],x=1..100);
\end{lstlisting}
\begin{Maple Normal}
\_\_\_\_\_\_\_\_\_\_\_\_\_\_\_\_\_\_\_\_\_\_\_\_\_\_\_\_\_\_\_\_\_\_\_\_\_\_\_\_\_\_\_\_\_\_\_\_\_\_\_\_\_\_\_\_\_\_\_\_\_\_\_\_\_\_\_\_\_\_\_\_\_\_\_\_\_\_\_\_\_
\end{Maple Normal}
\begin{Maple Heading 2}
\textbf{Exercise 1.6}
\end{Maple Heading 2}
\begin{lstlisting}
> g := (a,x) -> (x-1-sin(a/4))*(x-1-cos(a/4))*
              (x+1-sin(a/4))*(x+1-cos(a/4))/(1+x^4);
\end{lstlisting}
\begin{Maple Normal}
\textbf{(a)} The graph  
$y =g (2,x)$ approaches 
$y =1$ very closely when 
$x$ is very large (positive or negative).  Close to 
$x =0$ it climbs quickly to about 
$y = 2.5$, then drops very quickly to about 
$ y =0$, then climbs again to just below 
$y =1$.  The graphs for 
$g (3,x)$and 
$g (4,x)$ look very similar on this scale.
\end{Maple Normal}
\begin{lstlisting}
> plot(g(2,x),x=-500..500);
\end{lstlisting}
\begin{lstlisting}
> plot(g(3,x),x=-500..500);
\end{lstlisting}
\begin{lstlisting}
> plot([g(4,x),1],x=-500..500);
\end{lstlisting}
\begin{Maple Heading 2}
\textbf{(b)}
\end{Maple Heading 2}
\begin{lstlisting}
> plot(g(2,x),x=-1..2);
\end{lstlisting}
\begin{lstlisting}
> plot(g(4,x),x=-1..2);
\end{lstlisting}
\begin{lstlisting}
> plot(g(3,x),x=-1..2);
\end{lstlisting}
\begin{lstlisting}
> plot(g(3,x),x=-0.4..-0.2,-0.01..0,numpoints=1000);
\end{lstlisting}
\begin{Maple Normal}
It works out that the graph of 
$g (\pi,x)=g ( 3.1415927,x)$ stays above the 
$x$-axis, just touching it in two places.  We first show this graphically:
\end{Maple Normal}
\begin{lstlisting}
> plot(g(Pi,x),x=-1..3);
\end{lstlisting}
\begin{Maple Normal}
We now zoom in to the first point where the curve touches the axis:
\end{Maple Normal}
\begin{lstlisting}
> plot(g(Pi,x),x=-.292899..-.292885,-1e-10..1e-10,numpoints=1000);
\end{lstlisting}
\begin{Maple Normal}
To see why this works algebraically, note that 
$\sin (\frac{\pi}{4})=\frac{\sqrt{2}}{2}$ and 
$\cos (\frac{\pi}{4})=\frac{\sqrt{2}}{2}$.  Putting this in the definition of 
$g (a ,x)$ gives the following:
\end{Maple Normal}
\begin{lstlisting}
> g(Pi,x);
\end{lstlisting}
\begin{Maple Normal}
As the square or fourth power of any real number is nonnegative, this formula shows that 
$0\le g (\pi,x)$ for all 
$x$.
\end{Maple Normal}
\begin{Maple Normal}

\end{Maple Normal}
\begin{Maple Normal}
\textbf{(c) }It turns out that 
$1-8 g (a ,0)=\cos (a)$.  To see this graphically, compare the two plots below.
\end{Maple Normal}
\begin{lstlisting}
> plot(1-8*g(a,0),a=-10..10);
\end{lstlisting}
\begin{lstlisting}
> plot(cos(a),a=-10..10);
\end{lstlisting}
\begin{Maple Normal}
However, Maple's simplification commands are not clever enough to work this out:
\end{Maple Normal}
\begin{lstlisting}
> 1-8*g(a,0);
\end{lstlisting}
\begin{lstlisting}
> simplify(1-8*g(a,0));
\end{lstlisting}
\begin{Maple Normal}
If we give Maple a hint by asking it to simplify 
$1-8 g (a ,0)-\cos (a)$ instead, it does succeed in working out that this is zero:
\end{Maple Normal}
\begin{lstlisting}
> simplify(1-8*g(a,0)-cos(a));
\end{lstlisting}
\begin{Maple Normal}
Here is a trick that persuades Maple to convert 
$1-8 g (a ,0)$ to 
$\cos (a)$.  It is based on De Moivre's Theorem, which will be covered in the course on complex numbers.
\end{Maple Normal}
\begin{lstlisting}
> simplify(convert(1-8*g(a,0),exp));
\end{lstlisting}
\begin{Maple Normal}
To do this simplification by hand, note that 
$(-1-\cos (\frac{a}{4}))(1-\cos (\frac{a}{4}))=\cos (\frac{a}{4})^{2}-1$, which is the same as 
$-\sin (\frac{a}{4})^{2}$.  Similarly, we have 
$(-1-\sin (\frac{a}{4})) (1-\sin (\frac{a}{4}))=-\cos (\frac{a}{4})^{2}$.  Putting these together, we see that 
$8 g (a ,0)=8 \sin (\frac{a}{4})^{2} \cos (\frac{a}{4})^{2}$, which is the same as 
$2 (2 \sin (\frac{a}{4}) \cos (\frac{a}{4}))^{2}$.  Using the standard formula 
$\sin (2 \theta)=2 \sin (\theta) \cos (\theta)$, this simplifies to 
$2 \sin (\frac{a}{2})^{2}$.  We can then use the standard formula 
$\cos (2 \varphi)=1-2 \sin (\varphi)^{2}$ to deduce that 
$1-8 g (a ,0)=\cos (a)$. 
\end{Maple Normal}
\begin{Maple Heading 2}
\textbf{(d)}
\end{Maple Heading 2}
\begin{lstlisting}
> plot(g(a,0.5),a=-40..40);
\end{lstlisting}
\begin{lstlisting}
> slope := diff(g(a,0.5),a);
\end{lstlisting}
\begin{Maple Normal}
One of the tall peaks is at 
$x =\pi$.
\end{Maple Normal}
\begin{lstlisting}
> fsolve(slope=0,a=3);
\end{lstlisting}
\begin{Maple Normal}
The others are at 
$x =-7 \pi$ and 
$x =9 \pi$. 
\end{Maple Normal}
\begin{lstlisting}
> fsolve(slope=0,a=-22);
\end{lstlisting}
\begin{lstlisting}
> evalf(%/Pi);
\end{lstlisting}
\begin{lstlisting}
> fsolve(slope=0,a=29);
\end{lstlisting}
\begin{lstlisting}
> evalf(%/Pi);
\end{lstlisting}
\begin{Maple Heading 2}
\textbf{(e)}
\end{Maple Heading 2}
\begin{lstlisting}
> plot3d(g(a,x),a=-20..20,x=-15..15,
       axes=boxed,grid=[100,100]);
\end{lstlisting}
\begin{lstlisting}
> 
\end{lstlisting}
\begin{lstlisting}
> 
\end{lstlisting}
\begin{lstlisting}
> 
\end{lstlisting}
\end{document}
